\subsection{Uplifting Runaways --- arXiv:1809.06861}

\textbf{Authors:} Iosif Bena, Emilian Dudas, Mariana Gra\~na, Severin L\"ust

\paragraph{主な主張}
ワープした変形conifold内の反ブレーンが先端の球面サイズを不安定化し得ることを示し、その回避に必要な大きなフラックスが、解析対象のType IIBコンパクト化でtadpole条件およびスケール階層と緊張することを論じた。 \cite{Bena:2018fqc}

\paragraph{新規性}
喉部を固定背景として扱う近似を見直し、軽いconifold変形モジュライに沿うrunawayを反ブレーンupliftの制約として明示した。

\paragraph{位置付け}
反ブレーンによる真空構成で、喉部の安定化とコンパクト空間全体の電荷制約を同時に検証する必要性を示す研究である。

\paragraph{コメント（読書メモ）}
檜垣さんの科研費申請書の参考資料[3]。読書時には「KKLTで仮定されていた反ブレーンの効果を具体的に計算すると、conifold不安定化機構のため安定化が難しそうだ」と理解した。そこから、従来仮定していた効果を弦理論の立場で詳しく計算すれば、見落としていた効果が見つかるのではないかと考えた。
